% Figure: the delta (pi) network and its wye (tee) equivalent, drawn with
% plain TikZ paths and nodes (no circuitikz). The three terminals A, B, C are
% common to both drawings; the delta resistors sit on the edges of the
% triangle and the wye resistors meet at a central node N.
\begin{figure}[htbp]
\centering
\begin{tikzpicture}[
  line width=0.8pt,
  font=\scriptsize,
]
% ===== delta network (left) =====
\begin{scope}[shift={(0,0)}]
  \coordinate (A) at (0,3.2);
  \coordinate (B) at (-1.6,0);
  \coordinate (C) at (1.6,0);
  \node[anchor=south] at (A) {$A$};
  \node[anchor=east]  at (B) {$B$};
  \node[anchor=west]  at (C) {$C$};
  \filldraw (A) circle [radius=1.4pt];
  \filldraw (B) circle [radius=1.4pt];
  \filldraw (C) circle [radius=1.4pt];
  % edge A-B with resistor label R_AB
  \draw (A) -- (B);
  \node[anchor=east] at (-1.1,1.7) {$R_{AB}$};
  % edge A-C with resistor label R_AC
  \draw (A) -- (C);
  \node[anchor=west] at (1.1,1.7) {$R_{AC}$};
  % edge B-C with resistor label R_BC
  \draw (B) -- (C);
  \node[anchor=north] at (0,-0.05) {$R_{BC}$};
  \node[enggray, anchor=north] at (0,-0.6) {delta ($\Delta$)};
\end{scope}

% ===== arrow =====
\draw[->, engblue, line width=1.0pt] (2.7,1.4) -- (3.9,1.4);
\node[engblue, anchor=south, font=\tiny] at (3.3,1.5) {equivalent};

% ===== wye network (right) =====
\begin{scope}[shift={(6.6,0)}]
  \coordinate (A2) at (0,3.2);
  \coordinate (B2) at (-1.6,0);
  \coordinate (C2) at (1.6,0);
  \coordinate (N)  at (0,1.4);
  \node[anchor=south] at (A2) {$A$};
  \node[anchor=east]  at (B2) {$B$};
  \node[anchor=west]  at (C2) {$C$};
  \filldraw (A2) circle [radius=1.4pt];
  \filldraw (B2) circle [radius=1.4pt];
  \filldraw (C2) circle [radius=1.4pt];
  \filldraw[engorange] (N) circle [radius=1.6pt];
  \node[engorange, anchor=west] at (0.15,1.55) {$N$};
  % spokes
  \draw (A2) -- (N);
  \node[anchor=west] at (0.1,2.4) {$R_A$};
  \draw (B2) -- (N);
  \node[anchor=south east] at (-0.85,0.75) {$R_B$};
  \draw (C2) -- (N);
  \node[anchor=south west] at (0.85,0.75) {$R_C$};
  \node[enggray, anchor=north] at (0,-0.6) {wye ($Y$)};
\end{scope}

\end{tikzpicture}
\caption[Delta-wye transformation]{The three-terminal delta network and its
wye equivalent. The two boxes present identical terminal behaviour at $A$, $B$,
$C$ for every external circuit, so a bridge or ladder that resists
series-parallel reduction is unlocked by converting one delta to a wye (or the
reverse). Each wye spoke is the product of the two adjoining delta edges
divided by the sum of all three, $R_A = R_{AB}R_{AC}/(R_{AB}+R_{AC}+R_{BC})$,
with $R_B$ and $R_C$ following by rotating the labels.}
\label{fig:vol04ch09:delta-wye}
\end{figure}
